% Research fragment for integration into the parent article.
% All assertions in this fragment have a self-contained algebraic proof.
% The notation C_j,D_r,a_{jk} is the notation of the source tail-rank section.

\section{The optimal moment prefix for every pair \texorpdfstring{$(2,q)$}{(2,q)}}
\label{sec:rank-two}

The source conjecture asks whether the first $q$ centered moments have
full formal rank for every pair of even orders $2\leq p<q$.
We prove the whole infinite subfamily $p=2$, with an exact determinant,
an explicit rational extraction row, and an ordinary convergent
polygamma-tail identity.  The order-four and order-six families are proved below; unequal pairs $8\leq p<q$
remain open here.

The Pad\'e and Hankel ingredients used below are classical.  In
particular, the normalized shifted Bernoulli Hankel determinant is
recorded by Dilcher and Jiu~\cite[\S7.3, Eq.~(7.12)]{DilcherJiu2020},
who credit the evaluation to Krattenthaler.  We reprove the needed
formula to make the normalization and the inverse row explicit.
The new step relative to the source tail-moment report is the
rank-one perturbation calculation giving the factor $2-B_{q-2}$,
and its consequences for the optimal prefix and rational tail sums.

Write $q=2b$, where $b\geq2$, and put
\[
 M_b=(a_{b+t,k})_{0\leq t,k<b}.
\]
The entry convention is exactly the source convention: a binomial
summand is omitted if its lower index is less than $2$.  Thus, when
$p=2$, the first binomial summand occurs only in column $k=0$.

\begin{theorem}[Exact determinant for the entire order-two family]
\label{thm:rank-two}
For every integer $b\geq2$,
\begin{equation}\label{eq:rank-two-det}
 \det M_b=
 \frac{(-1)^{b-1}\bigl(2-B_{2b-2}\bigr)}
 {2^b(2b-2)!}\ne0.
\end{equation}
Consequently, for $p=2$ and $q=2b$, the first $q$ centered moments
$D_0,\ldots,D_{q-1}$ have full rank $q$ on the source's formal
coordinate list
\[
 C_0,\ldots,C_{b-1},\quad Z_0,\ldots,Z_{b-1},
 \qquad Z_k=\zeta(2k),\quad Z_0=-\tfrac12.
\]
The determinant of their full coefficient matrix, in this coordinate
order, is
\begin{equation}\label{eq:rank-two-full-det}
 \frac{(-1)^{b-1}\bigl(2-B_{2b-2}\bigr)}
 {2^b(2b-2)!\,(4b-1)!!}.
\end{equation}
This is a formal-rank assertion; it does not assert arithmetic
independence of the numerical coordinates.
\end{theorem}

\subsection{A finite Pad\'e calculation, proved explicitly}

We use Bernoulli numbers with generating function
\[
 \frac{t}{e^t-1}=\sum_{r\geq0}B_r\frac{t^r}{r!},
 \qquad B_1=-\frac12,
\]
and write
\[
 \mu_r=\frac{B_{2r}}{(2r)!},\qquad
 F(z)=\sum_{r\geq0}\mu_rz^r
     =\frac{\sqrt z}{2}\coth\frac{\sqrt z}{2}.
\]
The expression involving $\sqrt z$ is an even power series and hence
defines a single-valued holomorphic germ at $z=0$.

\begin{lemma}[Explicit exponential Pad\'e identity]
\label{lem:exp-pade-rank}
For $N\geq0$, define
\[
 P_N(t)=\sum_{k=0}^N
   \frac{(2N-k)!\,N!}{(2N)!\,k!\,(N-k)!}\,t^k.
\]
Then
\begin{equation}\label{eq:exp-pade-rank}
 e^tP_N(-t)-P_N(t)
 =\frac{(-1)^N(N!)^2}{(2N)!(2N+1)!}\,t^{2N+1}
  +O(t^{2N+2}).
\end{equation}
\end{lemma}

\begin{proof}
For $0\leq r\leq2N$, the coefficient of $t^r$ in
$e^tP_N(-t)$ equals
\[
 \frac{(2N-r)!}{(2N)!}
 \sum_{k=0}^N(-1)^k\binom Nk
                 \binom{2N-k}{2N-r}.
\]
The sum is the coefficient of $x^{2N-r}$ in
\[
 \sum_{k=0}^N(-1)^k\binom Nk(1+x)^{2N-k}
 =x^N(1+x)^N.
\]
It is therefore $\binom Nr$ for $r\leq N$ and zero for $N<r\leq2N$.
These are precisely the coefficients of $P_N(t)$ through degree
$2N$.  The next coefficient is
\[
 \frac1{(2N)!}\sum_{k=0}^N
       \frac{(-1)^k\binom Nk}{2N+1-k}
 =\frac1{(2N)!}\int_0^1x^N(x-1)^N\,dx,
\]
which gives the stated beta-integral value.
\end{proof}

For $n\geq0$, take $N=2n+1$ and form the polynomials
\begin{align}
 q_n(t^2)&=\frac{P_{2n+1}(t)-P_{2n+1}(-t)}t,
 \\
 r_n(t^2)&=\frac{P_{2n+1}(t)+P_{2n+1}(-t)}2.
\end{align}
Both have degree $n$ and $q_n(0)=1$.  Direct simplification gives
\begin{equation}\label{eq:pade-coth-rank}
 F(z)q_n(z)-r_n(z)
 =h_nz^{2n+1}+O(z^{2n+2}),
 \qquad
 h_n=\frac{((2n+1)!)^2}{(4n+2)!(4n+3)!}.
\end{equation}
Indeed, after setting $z=t^2$, the left side is exactly
\[
 -\frac{e^tP_{2n+1}(-t)-P_{2n+1}(t)}{e^t-1}.
\]
This proves both the vanishing coefficients and the leading
coefficient in \eqref{eq:pade-coth-rank}; no normality theorem for
Pad\'e approximants is being assumed.

Write
\begin{equation}\label{eq:q-explicit-rank}
 q_n(z)=\sum_{\ell=0}^nq_{n,\ell}z^\ell,
 \qquad
 q_{n,\ell}=
 \frac{2(4n+1-2\ell)!\,(2n+1)!}
 {(4n+2)!\,(2\ell+1)!\,(2n-2\ell)!},
\end{equation}
and reverse its coefficients:
\[
 \pi_n(x)=x^nq_n(1/x).
\]
This is a monic polynomial of degree $n$.
Let $\mathcal L$ be the linear functional on polynomials specified by
$\mathcal L(x^r)=\mu_r$.  Comparing coefficients in
\eqref{eq:pade-coth-rank} proves
\begin{align}
 \mathcal L(x^k\pi_n(x))&=0 &&(1\leq k\leq n),
       \label{eq:orthogonality-rank}\\
 \mathcal L(\pi_n)&=c_n,
 &c_n&=\frac{(2n+1)(2n+2)!}{(4n+2)!},
       \label{eq:constant-moment-rank}\\
 \mathcal L(x^{n+1}\pi_n)&=h_n.
       \label{eq:shifted-norm-rank}
\end{align}
Here \eqref{eq:constant-moment-rank} is the coefficient of $z^n$
in $r_n(z)$.

\subsection{The normalized Bernoulli Hankel determinant}

Put
\[
 H_b=(\mu_{i+j})_{0\leq i,j<b},\qquad
 \Delta_n=\det(\mu_{i+j+1})_{0\leq i,j<n},\qquad\Delta_0=1.
\]
The functional $\mathcal N(g)=\mathcal L(xg)$ makes the monic
polynomials $\pi_0,\pi_1,\ldots$ orthogonal by
\eqref{eq:orthogonality-rank}.  Their squared norms are $h_j$:
the polynomial $\pi_j(x)-x^j$ has degree less than $j$, so
\[
 \mathcal N(\pi_j^2)=\mathcal L(x^{j+1}\pi_j)=h_j.
\]
The triangular change from monomials to these monic polynomials has
determinant $1$.  Therefore
\begin{equation}\label{eq:delta-product-rank}
 \Delta_n=\prod_{j=0}^{n-1}
       \frac{((2j+1)!)^2}{(4j+2)!(4j+3)!}.
\end{equation}
In $H_b$, replace its last row by the linear combination of rows
whose coefficients are those of $\pi_{b-1}$.  This operation leaves
the determinant unchanged because $\pi_{b-1}$ is monic.
The resulting last row is $(c_{b-1},0,\ldots,0)$, whence
\begin{equation}\label{eq:H-product-rank}
 \det H_b=(-1)^{b-1}c_{b-1}\Delta_{b-1}.
\end{equation}
In particular, the Bernoulli Hankel matrix is nonsingular.

There is also a useful analytic explanation of this nonvanishing.
Euler's even-zeta identity gives, for $r\geq1$,
\[
 (-1)^{r-1}\mu_r
 =2\sum_{m\geq1}(2\pi m)^{-2r}.
\]
After conjugation by $\operatorname{diag}(1,-1,1,\ldots)$, the
bottom-right $(b-1)$-square block of $H_b$ is the negative of a
positive definite moment matrix.  Positivity is strict because a
nonzero polynomial cannot vanish at the infinitely many distinct
points $(2\pi m)^{-2}$.  The top-left entry remains $1$, so its
Schur complement is positive.  This yields the sign
$(-1)^{b-1}$ and invertibility independently of the exact product.
The exceptional initial coefficient $\mu_0=1$ is retained; replacing
it by a positive Stieltjes mass would change the problem.

\subsection{The mixed matrix is a rank-one perturbation}

The source coefficient formula, specialized to $p=2$, is
\begin{equation}\label{eq:rank-two-entry}
 (M_b)_{tk}
 =\frac{(2b+2t-1)!}{(2b-2k)!}\,\mu_{t+k}
   +\mathbf1_{k=0}\frac{2t+1}{2}B_{2t}.
\end{equation}
Define diagonal matrices
\[
 R_b=\operatorname{diag}_{0\leq t<b}(2b+2t-1)!,\qquad
 E_b=\operatorname{diag}_{0\leq k<b}\frac1{(2b-2k)!},
\]
and a column vector $u=(u_t)_{0\leq t<b}$ by
\[
 u_t=\mu_t\rho_t,\qquad
 \rho_t=\frac{(2b)!\,(2t+1)!}{2(2b+2t-1)!}.
\]
Then
\begin{equation}\label{eq:rank-one-factor}
 M_b=R_b(H_b+u e_0^{\mathsf T})E_b,
\end{equation}
where $e_0=(1,0,\ldots,0)^{\mathsf T}$.

Set $n=b-1$.  By \eqref{eq:orthogonality-rank} and
\eqref{eq:constant-moment-rank}, the first row of $H_b^{-1}$ is the
coefficient row of $\pi_n/c_n$.  Consequently
\begin{equation}\label{eq:contraction-rank}
 e_0^{\mathsf T}H_b^{-1}u
 =\frac1{c_{b-1}}\sum_{t=0}^{b-1}q_{b-1,b-1-t}\,\mu_t\rho_t.
\end{equation}
The factorials in \eqref{eq:q-explicit-rank} cancel to give
\begin{align}
 \frac1{c_{b-1}}\sum_{t=0}^{b-1}
       q_{b-1,b-1-t}\mu_t\rho_t
 &=\frac1{2b-1}\sum_{t=0}^{b-1}
       \binom{2b-1}{2t}(2t+1)B_{2t}\\
 &=1-B_{2b-2}.
       \label{eq:bernoulli-contraction-rank}
\end{align}
For clarity, the last equality is an elementary Bernoulli identity.
If $N=2b-1\geq3$, then
\[
 \sum_{k=0}^N\binom NkB_k=B_N(1)=0,
 \qquad
 \sum_{k=0}^Nk\binom NkB_k=-NB_{N-1}.
\]
Removing the $k=1$ terms and using $B_k=0$ for odd $k>1$ shows
\[
 \sum_{\substack{0\leq k<N\\ k\text{ even}}}
      (k+1)\binom NkB_k=N(1-B_{N-1}),
\]
as required.  The condition $b\geq2$ matters: this identity is not
the corresponding formula for $N=1$.

The matrix determinant lemma now yields
\[
 \det(H_b+u e_0^{\mathsf T})
 =\det H_b\bigl(2-B_{2b-2}\bigr).
\]
It remains to simplify the factorial factor.  Equations
\eqref{eq:delta-product-rank}--\eqref{eq:H-product-rank} give
\begin{equation}\label{eq:factorial-cancel-rank}
 \det R_b\det E_b\det H_b
 =\frac{(-1)^{b-1}}{2^b(2b-2)!}.
\end{equation}
One quick verification, also a proof for every $b$, is induction:
the left side is $1/2$ at $b=1$ and the quotient of its values at
$b$ and $b-1$ is
\[
 -\frac1{4(b-1)(2b-3)}.
\]
This is exactly the quotient on the right side of
\eqref{eq:factorial-cancel-rank}.

For completeness, $B_{2m}$ cannot equal $2$ for any $m\geq1$ by an
elementary parity argument; the full von Staudt--Clausen theorem is
unnecessary.  Work in the ring $\mathbb Z_{(2)}$ of rational numbers
with odd denominator.  The Bernoulli recurrence gives
\[
 (2m+1)(2B_{2m})
 =2m-1-\sum_{j=1}^{m-1}\binom{2m+1}{2j}(2B_{2j}).
\]
Inductively $2B_{2j}\in\mathbb Z_{(2)}$ and is congruent to $1$
modulo $2$.  The binomial sum in the last display is even modulo
$2$, since
\[
 \sum_{j=1}^{m-1}\binom{2m+1}{2j}
 =2^{2m}-2m-2.
\]
Thus $2B_{2m}$ is also congruent to $1$ modulo $2$.  In particular
its $2$-adic valuation is $0$, whereas $2\cdot2$ is even.
This proves $2-B_{2m}\ne0$ and completes
\eqref{eq:rank-two-det}.  Finally the forward Bernoulli transform
has diagonal entries $1/(2r+1)$, giving
\eqref{eq:rank-two-full-det}.

\subsection{An all-order rational extraction identity}

The same calculation identifies one entire row of the inverse,
without an unspecified matrix inversion.

\begin{theorem}[Explicit extraction of the rational coordinate]
\label{thm:rational-row-two}
For $b\geq2$,
\begin{equation}\label{eq:rational-block-two}
 \sum_{t=0}^{b-1}\binom{2b-1}{2t}\,C_{b+t}
 =-\frac{(2b-1)(2-B_{2b-2})}{4}.
\end{equation}
Equivalently, with
\begin{equation}\label{eq:moment-weights-two}
 w_{b,r}=\sum_{\substack{0\leq t<b\\r\leq b+t}}
 \binom{2b-1}{2t}\binom{2b+2t+1}{2r}
       2^{-2(b+t-r)},\qquad 0\leq r\leq2b-1,
\end{equation}
the centered moments satisfy
\begin{equation}\label{eq:rational-moment-two}
 \sum_{r=0}^{2b-1}w_{b,r}D_r
 =-\frac{(2b-1)(2-B_{2b-2})}{4}.
\end{equation}
\end{theorem}

\begin{proof}
The rank-one inverse formula and
\eqref{eq:bernoulli-contraction-rank} show that
\[
 e_0^{\mathsf T}(H_b+u e_0^{\mathsf T})^{-1}
 =\frac{e_0^{\mathsf T}H_b^{-1}}{2-B_{2b-2}}.
\]
Insert the coefficient formula for $q_{b-1}$ and the diagonal
factors in \eqref{eq:rank-one-factor}.  The resulting first row is
\begin{equation}\label{eq:inverse-first-row-two}
 (M_b^{-1})_{0t}=
 \frac{2}{(2b-1)(2-B_{2b-2})}\binom{2b-1}{2t}.
\end{equation}
Multiplication by $(C_b,\ldots,C_{2b-1})^{\mathsf T}$ recovers
$Z_0=-1/2$, which proves \eqref{eq:rational-block-two}.
Substitution of the source inverse Bernoulli transform
\[
 C_j=\sum_{r=0}^j\binom{2j+1}{2r}2^{-2(j-r)}D_r
\]
gives \eqref{eq:rational-moment-two}.
\end{proof}

As an exact check, for $b=3$ the block identity is
\[
 C_3+10C_4+5C_5=-\frac{61}{24}.
\]
Multiplying its moment form by $1024$ reproduces the source's
$(2,6)$ identity with coefficients
$61,8204,116000,370048,303360,56320$ and right side $-7808/3$.
The source example is therefore a specialization of an all-order
formula, rather than an isolated elimination.

\subsection{An ordinary convergent series at every even order}

The finite-part identity has a direct, absolutely convergent
interpretation with an explicitly specified polynomial subtraction.
Set
\begin{align}
 S_b(y)&=\frac{y^{2b+1}}2
       \bigl((1+y)^{2b-1}+(1-y)^{2b-1}\bigr),
       \label{eq:S-polynomial-two}\\
 W_b(x)&=S_b(x+1/2)-S_b(x-1/2)
        =\sum_{r=0}^{2b-1}w_{b,r}x^{2r}.
       \label{eq:W-polynomial-two}
\end{align}
For an even integer $s\geq2$, define rational coefficients
\[
 c_{s,0}=\frac1{s-1},\qquad
 c_{s,h}=\frac{B_{2h}(1/2)}{(2h)!}(s)_{2h-1}\quad(h\geq1),
\]
where $(s)_r$ is a rising factorial.  Put
\[
 T_{b,h}=\sum_{i=0}^hc_{2,i}c_{2b,h-i},
\]
and define the explicit even polynomial
\begin{equation}\label{eq:subtraction-two}
 R_b(x)=\sum_{r=b}^{2b-1}\ \sum_{h=0}^{r-b}
       w_{b,r}T_{b,h}x^{2(r-b-h)}.
\end{equation}
Only the finitely many coefficients with $h\leq b-1$ occur.

\begin{corollary}[Convergent rational tail sums]
\label{cor:ordinary-two}
Let $x_n=n+1/2$ and $A_s(n)=\zeta(s)-H_n^{(s)}$.  For every
$b\geq2$,
\begin{equation}\label{eq:ordinary-two}
 \sum_{n=0}^{\infty}
 \left\{W_b(x_n)A_2(n)A_{2b}(n)-R_b(x_n)\right\}
 =-\frac{(2b-1)(2-B_{2b-2})}{4}.
\end{equation}
The summand is $O(n^{-2})$, so the series is absolutely convergent.
Equivalently, one may replace the tail product by
\[
 A_2(n)A_{2b}(n)
 =\frac{\psi'(n+1)\psi^{(2b-1)}(n+1)}{(2b-1)!}.
\]
\end{corollary}

\begin{proof}
The centered Hurwitz expansion gives
\[
 A_s(n)\sim x_n^{1-s}\sum_{h\geq0}c_{s,h}x_n^{-2h},
 \qquad
 A_2(n)A_{2b}(n)\sim
       x_n^{-2b}\sum_{h\geq0}T_{b,h}x_n^{-2h}.
\]
The polynomial $R_b$ in \eqref{eq:subtraction-two} is precisely the
sum of all nonnegative powers in the product with $W_b$.
After its subtraction the next possible term is $O(x_n^{-2})$.
Every removed term has continued sum zero because
$\zeta(-2r,1/2)=0$ for $r\geq0$.  Thus the absolutely convergent
sum equals the finite-part moment combination in
\eqref{eq:rational-moment-two}.  Alternatively, the partial sum of
each removed centered even power is an odd polynomial in the upper
integer cutoff with no constant term, which proves the same
constant-term equality without interchanging a divergent series.
The polygamma expression is the standard positive-even-order
identity $A_s(n)=\psi^{(s-1)}(n+1)/(s-1)!$.
\end{proof}

For example, when $b=2$, multiplication by $64$ gives the compact
ordinary identity
\begin{equation}\label{eq:ordinary-24}
 \sum_{n=0}^{\infty}\left\{
 (7+412x_n^2+2000x_n^4+1344x_n^6)A_2(n)A_4(n)
 -448x_n^2-\frac{1216}{3}
 \right\}=-88.
\end{equation}

An independently reconstructed example at $b=4$ is also useful for
checking the size and sign of the subtraction. Define
\begin{align}
 \mathsf P_4(X)={}&547+129756X^2+3657584X^4+27896512X^6\notag\\
 &+71997696X^8+66311168X^{10}+20500480X^{12}+1720320X^{14},
 \label{eq:P28-example}\\
 \mathsf R_4(X)={}&245760X^6+2334720X^4+3955712X^2
                                    +\frac{37765120}{63}.
 \label{eq:R28-example}
\end{align}
These are $2^{14}W_4$ and $2^{14}R_4$, respectively. Then
\begin{equation}\label{eq:ordinary-28}
 \boxed{\displaystyle
 \sum_{n=0}^{\infty}
    \{\mathsf P_4(x_n)A_2(n)A_8(n)-\mathsf R_4(x_n)\}
       =-\frac{169984}{3}.}
\end{equation}
The associated exact coefficient matrix has determinant
$\det M_4=-83/483840$ and extraction weights $(1,21,35,7)$.
The numerical example is not used to infer the formula: every coefficient
above follows from the proved finite polynomials.

\paragraph{What this resolves and what remains open.}
The optimal consecutive-prefix conjecture is now proved for every
even pair with smaller order $2$, as well as the diagonal family
already settled in the source.  The determinant factor
$2-B_{q-2}$ explains why a uniform real-sign argument is misleading:
its sign can change as the even Bernoulli numbers grow.  The nonzero
factor is instead proved arithmetically.  The following sections extend the determinant method to $p=4,6$ and give
an exact smaller determinant for every even pair. Nonvanishing for all
unequal pairs $8\le p<q$ remains open.
